\textbf{Application: auditing and improving AI-assisted scientific review.}
A review can sound plausible or resemble human feedback while adding little to the information available to a reviewer or area chair.
My ongoing work audits LLM-generated reviews through complementary information value.
Preliminary analyses indicate that raw LLM reviews offer little complementary value, while structured pipelines can provide more.
This motivates evaluating review assistance conditional on existing expert judgments, and distinguishing the information a system supplies from how persuasively it presents that information.

In joint work on learning review checklists~\cite{liu2026checklists}, we develop an intervention that helps reviewers cover evidence they might otherwise miss.
Our LLM-assisted pipeline extracts atomic claims from historical reviews, retains concerns verifiable against manuscript materials, and organizes them into a hierarchy of checklist items.
We select the granularity using correlated agreement: agreement between independent reviews of the same paper beyond agreement across different papers.
Shared checks can reduce reviewer lottery variance but also introduce common bias.
In simulations on held-out NeurIPS 2021 and ICLR 2022 papers, the learned checklists reduced three-reviewer committee mean squared error by 58\% and 21\%, respectively, relative to a simulated unaided-reviewer population target.
In a 20-participant study, LLM-completed checklists reduced paper-ranking mean squared error from 0.48 to 0.22 against the NeurIPS ordering.
These targets assess consistency with reference judgments rather than objective scientific quality.

These projects connect diagnosing missing information to designing assistance that supplies it.
They motivate a broader question for AI+Science: how can AI expand the evidence experts consider while preserving their specialized knowledge and independent judgment?
